\chapter{Conclusões Parciais}
\label{cap:12}

Este trabalho propôs o desenvolvimento de um sistema especialista educacional para
validação de compatibilidade de \textit{hardware}, fundamentado no formalismo de Problemas de
Satisfação de Restrições e no algoritmo de propagação AC-3, com a finalidade de não apenas
detectar incompatibilidades entre componentes, mas de comunicar ao usuário leigo a razão
técnica de cada uma em linguagem acessível. Nesta etapa de Validação de Projeto,
apresentam-se as conclusões parciais obtidas a partir do que já foi elicitado, modelado e
implementado.

O objetivo geral mostrou-se plenamente viável ao longo do desenvolvimento. A modelagem do
problema de compatibilidade de \textit{hardware} como um CSP revelou-se adequada, uma vez que as
três categorias de restrições estudadas, soquete, padrão de memória e capacidade
energética, ajustaram-se naturalmente à estrutura de variáveis, domínios e restrições.
Sobre essa modelagem, o algoritmo AC-3 foi implementado como motor de inferência funcional,
capaz de propagar restrições de forma incremental e podar os domínios das variáveis a cada
seleção do usuário, sem recorrer à avaliação exaustiva por força bruta. Essa afirmação
deixou de ser apenas teórica: a avaliação de desempenho conduzida no
Capítulo~\ref{cap:resultados} comprovou empiricamente que a propagação de restrições decide
a consistência da rede em ordens de magnitude menos avaliações que a enumeração exaustiva,
a qual sequer conclui dentro do teto estabelecido a partir de domínios com cinquenta
componentes por categoria.

Quanto aos objetivos específicos, os resultados parciais são igualmente favoráveis. A
representação de conhecimento adotada, baseada em um modelo genérico de metadados, permitiu
dissociar o motor de inferência do domínio de aplicação, de modo que as regras de
compatibilidade residem inteiramente na base de conhecimento e não no código. Essa decisão
concretizou os requisitos de expansibilidade estabelecidos, pois a inclusão de novas
categorias, características ou restrições depende apenas da inserção de registros, sem
alteração do algoritmo. O mecanismo de explicação também foi implementado, gerando para
cada incompatibilidade uma justificativa em linguagem natural a partir de modelos de texto
associados às restrições, o que assegura que nenhum bloqueio ocorra de forma silenciosa.

A aplicação do motor ao domínio de compatibilidade de \textit{hardware}, tomado como estudo de caso,
confirmou a validade da abordagem. A base de conhecimento foi populada com componentes
reais das plataformas AM4 e AM5, e os cenários de teste demonstraram que o sistema
identifica corretamente as incompatibilidades de soquete, memória e potência, apresentando
as justificativas correspondentes. A interface no formato wizard, por sua vez, materializou
o caráter educativo da solução ao manter os componentes incompatíveis visíveis e
acompanhados de suas explicações, em vez de simplesmente ocultá-los.

Do ponto de vista do problema de pesquisa, os resultados parciais indicam que é possível
construir um sistema especialista que valide a compatibilidade de \textit{hardware} por meio de um
motor formal e que, ao mesmo tempo, ensine o usuário comunicando a razão de cada
incompatibilidade. A articulação entre a validação por propagação de restrições e a
capacidade de explicação herdada dos sistemas especialistas mostrou-se não apenas possível,
mas coerente, atendendo simultaneamente à dimensão técnica e à dimensão pedagógica que
motivaram o trabalho. A evidência experimental reforça essa conclusão sob dois regimes
complementares. Na ausência de seleção, a vantagem do motor é de natureza algorítmica, pois
a consistência é decidida sem enumerar o espaço de configurações completas. Sob seleção
parcial, situação que caracteriza o uso real do sistema, soma-se a esse ganho a redução
efetiva do espaço de busca por propagação, o que confirma a adequação da abordagem à
interação incremental do usuário com o wizard.

Reconhecem-se, contudo, as limitações próprias do estágio atual. A corretude das validações
depende do cadastro correto das restrições na base de conhecimento, característica inerente
aos sistemas dirigidos a dados. A validação da capacidade energética é realizada de forma
binária, componente a componente, preservando a natureza do grafo de restrições, ao passo
que o consumo agregado da configuração é apresentado apenas como informação. Além disso, a
estratégia de testes automatizados encontra-se parcialmente desenvolvida, com os testes de
integração e de interface previstos como atividade posterior.

Para a etapa de Avaliação Final, planeja-se a ampliação da base de conhecimento para
contemplar outras plataformas e categorias de componentes, a complementação da estratégia
de testes automatizados e a realização de uma avaliação da eficácia educativa da solução
junto a usuários. As conclusões parciais aqui apresentadas confirmam que os objetivos
propostos são alcançáveis e que o sistema desenvolvido cumpre, até o presente estágio, os
propósitos técnico e educacional que fundamentaram este trabalho.